\section{RED: Recurrent Event Detection}

\subsection{Notation}
For one class, let
\begin{itemize}
\item $e_t\in\{0,1\}$ be the event-presence variable at frame $t$;
\item $s_t=P(e_t=1\mid e_{t-1}=0)$ be the conditional start probability;
\item $z_t=P(e_t=0\mid e_{t-1}=1)$ be the conditional end probability;
\item $p_t=P(e_t=1)$ be the presence probability.
\end{itemize}
The model produces start and end logits; a sigmoid gives $s_t$ and $z_t$.

\subsection{The recurrence}
Assume $p_{-1}=0$. Presence at $t$ can come from two cases: an event starts at
$t$ (probability $s_t$, previous frame inactive with probability $1-p_{t-1}$),
or an event was active and does not end at $t$ (probability $1-z_t$, previous
frame active with probability $p_{t-1}$). Therefore
\begin{equation}
\boxed{p_t = s_t(1-p_{t-1}) + (1-z_t)p_{t-1}}.
\end{equation}
The onset and offset probabilities are
\begin{equation}
\mathrm{onset}_t = s_t(1-p_{t-1}),
\qquad
\mathrm{offset}_t = z_t p_{t-1}.
\end{equation}
Presence rises at onsets and falls at offsets, while remaining smooth
elsewhere.

\subsection{Why it is cheap}
Rewrite the recurrence as an affine scan:
\begin{equation}
p_t = a_t p_{t-1} + b_t,
\qquad
a_t = 1-s_t-z_t,
\qquad
b_t = s_t.
\end{equation}
Affine recurrences compose associatively:
\begin{equation}
(a_1,b_1)\circ(a_2,b_2)
=
(a_2a_1,\; a_2b_1+b_2).
\end{equation}
This permits a parallel prefix scan in $O(\log T)$ sequential steps instead of
$T$ Python iterations. The implementation uses a Hillis--Steele scan \cite{heinsen2023}; a sequential
reference is kept for tests. RED has no learnable parameters.

\subsection{Why this helps}
The acoustic model no longer learns one presence curve directly. It learns two
simpler conditional ``switch'' signals: when an event starts and when it ends.
Presence is derived deterministically. The onset and offset outputs are
temporally localised peaks, exactly what the boundary losses can supervise.

